% =====================================================================
% pertemuan-10.tex
% Pertemuan 10 - Bottom-Up Parsing: Shift-Reduce dan Operator Precedence
% Sub-CPMK 4.2 - Praktik terbimbing
% =====================================================================
\documentclass[aspectratio=169,10pt,t]{beamer}
\usepackage{config/beamerthemeubb}

\UBBmeeting{10}

\begin{document}

\UBBtitleframe{10}{Bottom-Up Parsing: Shift-Reduce dan Operator Precedence}{Praktik terbimbing}

% ---------------------------------------------------------------------
\begin{frame}{Capaian Pertemuan}
  \textbf{CPMK induk:} CPMK 04 --- Mengimplementasikan parser sederhana
  untuk validasi kode.

  \vspace{3pt}
  \begin{block}{Sub-CPMK yang dinilai}
    \small
    \textbf{Sub-CPMK 4.2:} Membuat tabel aksi dan \textit{goto} untuk
    parser SLR/LALR, lalu menguji dan men-\textit{debug} konflik
    \textit{shift-reduce} atau \textit{reduce-reduce} yang muncul.
  \end{block}

  \vspace{3pt}
  \begin{exampleblock}{Cakupan dua pertemuan}
    \footnotesize
    Pertemuan 10 menyiapkan model \eng{shift-reduce} dan tabel
    operator precedence. Tabel ACTION/GOTO SLR/LALR diselesaikan pada
    Pertemuan 11; keduanya menjadi artefak Sub-CPMK 4.2.
  \end{exampleblock}
\end{frame}

% ---------------------------------------------------------------------
\begin{frame}{Agenda Pertemuan (150 menit)}
  \small
  \begin{tabular}{@{}p{1.4cm}p{12.3cm}@{}}
    \toprule
    \textbf{Menit} & \textbf{Kegiatan} \\
    \midrule
    30 & Ceramah: model \textit{shift-reduce}, stack, dan handle \\
    30 & Latihan: menelusuri stack dan menandai handle \\
    40 & Lokakarya: tabel operator precedence \\
    30 & Praktik: kalkulator \textit{shift-reduce} dalam C \\
    20 & Diskusi konflik dan jembatan ke LR \\
    \bottomrule
  \end{tabular}

  \vspace{6pt}
  {\scriptsize\color{ubbblue!75}Porsi praktik dan latihan 100 menit.
  Keluaran hari ini menjadi data pendukung tabel ACTION/GOTO pada
  Pertemuan 11.}
\end{frame}

% ---------------------------------------------------------------------
\begin{frame}{Posisi Pertemuan pada Rantai Proyek}
  \UBBchain{3}

  \vspace{2pt}
  \begin{itemize}
    \item Pertemuan ini memperdalam simpul \textbf{parser} dengan arah
          bottom-up: dari daun menuju akar.
    \item Model \eng{shift-reduce} dan tabel precedence menjadi
          fondasi tabel ACTION/GOTO SLR/LALR.
    \item Konflik yang ditemukan hari ini didokumentasikan sebagai
          bahan \eng{debug} pada Pertemuan 11.
  \end{itemize}
\end{frame}

% ---------------------------------------------------------------------
\begin{frame}{Model Shift-Reduce}
  \begin{center}
    \begin{tikzpicture}[
        box/.style={draw=ubbblue!40,fill=ubbblue!7,rounded corners=3pt,
          minimum width=3.4cm,minimum height=1.5cm,align=center,
          font=\scriptsize,text=ubbblue},
        act/.style={box,draw=ubbgold,fill=ubbgold!20,minimum width=2.8cm}
      ]
      \node[box] (stack) {Stack\\[-1pt]\(\vdots\)\\handle};
      \node[box,right=0.5cm of stack] (input)
        {Masukan\\[-1pt]\(\code{id}+\code{id}*\code{id}\ \$\)};
      \node[act,right=0.5cm of input] (action) {Aksi\\shift / reduce};
      \draw[-{Stealth[length=1.6mm]},ubbblue!55,thick] (stack) -- (action);
      \draw[-{Stealth[length=1.6mm]},ubbblue!55,thick] (input) -- (action);
    \end{tikzpicture}
  \end{center}

  \begin{itemize}
    \item Parser membaca token kiri ke kanan dan memutuskan
          \eng{shift} atau \eng{reduce}.
    \item \eng{Shift}: dorong token ke stack.
    \item \eng{Reduce}: ganti ruas kanan aturan dengan nonterminal.
    \item Terima ketika stack berisi simbol awal dan masukan habis.
  \end{itemize}
\end{frame}

% ---------------------------------------------------------------------
\begin{frame}{Tata Bahasa Contoh dan Handle}
  \begin{block}{Tata bahasa ekspresi}
    \small
    \(E \to E+T \mid T\),\quad \(T \to T*F \mid F\),\quad
    \(F \to \code{id}\).
  \end{block}

  \begin{itemize}
    \item Tata bahasa ini \textit{left-recursive}, sehingga cocok untuk
          bottom-up tetapi tidak cocok untuk top-down tanpa
          transformasi.
    \item \textbf{Handle}: substring yang cocok dengan ruas kanan
          aturan dan dapat direduksi menjadi nonterminal kiri.
    \item Pada \(\code{id}+\code{id}*\code{id}\), handle pertama adalah
          \code{id} paling kiri.
    \item \code{id} direduksi menjadi \(F\), lalu \(T\), lalu \(E\).
  \end{itemize}
\end{frame}

% ---------------------------------------------------------------------
\begin{frame}{Penelusuran Shift-Reduce}
  \scriptsize
  \setlength{\tabcolsep}{4pt}
  \begin{tabular}{@{}p{2.8cm}p{4.4cm}p{5.8cm}@{}}
    \toprule
    \textbf{Stack} & \textbf{Masukan} & \textbf{Aksi} \\
    \midrule
     & \(\code{id}+\code{id}*\code{id}\ \$\) & shift \\
    \(\code{id}\) & \(+\code{id}*\code{id}\ \$\) & reduce \(F \to \code{id}\) \\
    \(F\) & \(+\code{id}*\code{id}\ \$\) & reduce \(T \to F\) \\
    \(T\) & \(+\code{id}*\code{id}\ \$\) & reduce \(E \to T\) \\
    \(E\) & \(+\code{id}*\code{id}\ \$\) & shift \(+\) \\
    \(E+\) & \(\code{id}*\code{id}\ \$\) & shift \code{id} \\
    \(E+\code{id}\) & \(*\code{id}\ \$\) & reduce \(F \to \code{id}\) \\
    \(E+F\) & \(*\code{id}\ \$\) & reduce \(T \to F\) \\
    \(E+T\) & \(*\code{id}\ \$\) & shift \(*\) \\
    \multicolumn{3}{@{}l}{\scriptsize\color{ubbblue!70}\ldots{} reduce
      \(F \to \code{id}\), \(T \to T*F\), \(E \to E+T\), lalu terima} \\
    \(E\) & \(\$\) & terima \\
    \bottomrule
  \end{tabular}

  \vspace{3pt}
  {\scriptsize\color{ubbblue!75}Keputusan diambil dari isi stack dan
  token saat ini, bukan dari seluruh masukan.}
\end{frame}

% ---------------------------------------------------------------------
\begin{frame}{Keuntungan dan Keterbatasan Bottom-Up}
  \begin{columns}[T,onlytextwidth]
    \begin{column}{0.48\textwidth}
      \begin{block}{Keuntungan}
        \footnotesize
        \begin{itemize}
          \item menangani \textit{left recursion} secara alami;
          \item mendukung grammar yang lebih luas;
          \item menjadi dasar parser LR;
          \item cocok untuk parser generator.
        \end{itemize}
      \end{block}
    \end{column}
    \begin{column}{0.48\textwidth}
      \begin{block}{Keterbatasan}
        \footnotesize
        \begin{itemize}
          \item tabel dan algoritma lebih kompleks;
          \item pesan kesalahan lebih sulit dibuat kontekstual;
          \item memerlukan konstruksi state yang sistematis.
        \end{itemize}
      \end{block}
    \end{column}
  \end{columns}

  \vspace{4pt}
  {\footnotesize\color{ubbblue!80}Top-down dan bottom-up bukan
  pengganti satu sama lain; pilihan bergantung pada grammar dan
  kebutuhan diagnostik.}
\end{frame}

% ---------------------------------------------------------------------
\begin{frame}{Viable Prefix}
  \begin{block}{Definisi}
    \small
    \textit{Viable prefix} adalah prefiks dari bentuk kalimat yang
    dapat muncul pada stack parser \textit{shift-reduce} tanpa melewati
    handle.
  \end{block}

  \begin{itemize}
    \item Menentukan apakah stack masih dapat diperluas menjadi bentuk
          kalimat valid.
    \item Jika stack bukan \textit{viable prefix}, parser tidak mungkin
          mencapai simbol awal.
    \item Viable prefix menjadi dasar pembentukan state LR(0) pada
          Pertemuan 11.
  \end{itemize}
\end{frame}

% ---------------------------------------------------------------------
\begin{frame}{Konflik Shift-Reduce}
  \begin{block}{Definisi}
    \small
    Konflik \textit{shift-reduce} terjadi ketika parser dapat memilih
    antara shift dan reduce untuk stack dan token yang sama.
  \end{block}

  \begin{exampleblock}{Contoh}
    \footnotesize
    Pada \(E \to E+E \mid E*E \mid \code{id}\), setelah stack berisi
    \(E+E\) dan token berikutnya \(*\), parser dapat:
    \begin{itemize}
      \item reduce \(E+E\) menjadi \(E\), atau
      \item shift \(*\) untuk membentuk \(E*E\).
    \end{itemize}
    Pilihan ini menentukan apakah penjumlahan atau perkalian
    dikelompokkan lebih dulu.
  \end{exampleblock}
\end{frame}

% ---------------------------------------------------------------------
\begin{frame}{Konflik Reduce-Reduce dan Penyelesaiannya}
  \begin{block}{Konflik \textit{reduce-reduce}}
    \small
    Terjadi ketika dua aturan berbeda dapat diterapkan pada handle yang
    sama. Biasanya menandakan grammar terlalu ambigu atau aturan
    tumpang tindih.
  \end{block}

  \begin{itemize}
    \item Selesaikan dengan mengubah grammar: precedence,
          associativity, atau factoring.
    \item Tambahkan aturan prioritas pada parser generator.
    \item Pilih shift atau reduce dengan kebijakan tetap yang
          didokumentasikan.
    \item Bison secara default memilih shift; keputusan ini dapat
          menyembunyikan bug jika tidak disengaja.
  \end{itemize}
\end{frame}

% ---------------------------------------------------------------------
\begin{frame}{Operator Precedence Grammar}
  \begin{block}{Syarat}
    \small
    \begin{enumerate}
      \item tidak memiliki produksi \(\varepsilon\);
      \item tidak memiliki dua nonterminal bersebelahan pada ruas
            kanan.
    \end{enumerate}
  \end{block}

  \begin{itemize}
    \item Parser dapat menentukan hubungan precedence antara terminal
          berdekatan tanpa tabel LR lengkap.
    \item Tiga relasi: \(a <\cdot b\) (lebih rendah),
          \(a \doteq b\) (sama), \(a \cdot> b\) (lebih tinggi).
    \item Simbol \(<\), \(=\), \(>\) di sini adalah relasi precedence,
          bukan operator aritmetika.
  \end{itemize}
\end{frame}

% ---------------------------------------------------------------------
\begin{frame}{Tabel Operator Precedence}
  \scriptsize
  \setlength{\tabcolsep}{5pt}
  \begin{tabular}{@{}c c c c c c@{}}
    \toprule
     & \texttt{id} & \texttt{+} & \texttt{*} & \texttt{(} &
       \texttt{)} \\
    \midrule
    \texttt{id} & & \(\cdot>\) & \(\cdot>\) & & \(\cdot>\) \\
    \texttt{+} & \(<\cdot\) & \(\cdot>\) & \(<\cdot\) & \(<\cdot\) &
      \(\cdot>\) \\
    \texttt{*} & \(<\cdot\) & \(\cdot>\) & \(\cdot>\) & \(<\cdot\) &
      \(\cdot>\) \\
    \texttt{(} & \(<\cdot\) & \(<\cdot\) & \(<\cdot\) & \(<\cdot\) &
      \(\doteq\) \\
    \texttt{)} & & \(\cdot>\) & \(\cdot>\) & & \(\cdot>\) \\
    \bottomrule
  \end{tabular}

  \vspace{5pt}
  {\scriptsize\color{ubbblue!75}Sel kosong menandakan kesalahan
  sintaksis. Tabel ini lebih kecil daripada tabel LR karena hanya
  memuat terminal, bukan himpunan item.}
\end{frame}

% ---------------------------------------------------------------------
\begin{frame}{Contoh Relasi Precedence}
  \small
  Untuk \(\code{id}+\code{id}*\code{id}\):
  \begin{itemize}
    \item \(\code{id} \cdot> +\): \code{id} harus direduksi sebelum
          \(+\).
    \item \(+ <\cdot \code{id}\): \(+\) menunggu operand.
    \item \(+ <\cdot *\): \(*\) memiliki precedence lebih tinggi.
    \item \(* \cdot> \code{id}\): \(*\) menunggu operand.
  \end{itemize}

  \vspace{4pt}
  \begin{block}{Cara membaca}
    \footnotesize
    Parser membandingkan terminal teratas stack dengan token saat ini.
    Relasi menentukan shift atau reduce: \(<\cdot\) dan \(\doteq\)
    mendorong shift; \(\cdot>\) memicu reduce.
  \end{block}
\end{frame}

% ---------------------------------------------------------------------
\begin{frame}{Algoritma Operator Precedence}
  \begin{enumerate}
    \setlength{\itemsep}{2pt}
    \item Cari terminal teratas pada stack, sebut \(a\).
    \item Bandingkan \(a\) dengan token saat ini \(b\).
    \item Jika \(a <\cdot b\), shift \(b\).
    \item Jika \(a \doteq b\), shift atau tangani kurung.
    \item Jika \(a \cdot> b\), reduce handle teratas.
    \item Jika sel kosong, laporkan kesalahan.
  \end{enumerate}

  \vspace{4pt}
  {\footnotesize Handle adalah substring stack yang cocok dengan ruas
  kanan aturan. Setelah reduce, bandingkan kembali terminal teratas
  dengan token saat ini.}
\end{frame}

% ---------------------------------------------------------------------
\begin{frame}{Contoh Penelusuran Operator Precedence}
  \small
  Untuk \(\code{id}+\code{id}*\code{id}\):
  \begin{enumerate}
    \setlength{\itemsep}{1pt}
    \item shift \code{id}, reduce ke \(F\), lalu \(T\), lalu \(E\).
    \item shift \(+\).
    \item shift \code{id}, reduce ke \(F\), lalu \(T\).
    \item \(+ <\cdot *\), shift \(*\).
    \item shift \code{id}, reduce ke \(F\).
    \item \(* \cdot> \$\): reduce \(T \to T*F\).
    \item \(+ \cdot> \$\): reduce \(E \to E+T\).
    \item terima.
  \end{enumerate}

  \vspace{3pt}
  {\footnotesize\color{ubbblue!80}Urutan ini memaksa perkalian
  direduksi sebelum penjumlahan.}
\end{frame}

% ---------------------------------------------------------------------
\begin{frame}[fragile]{Kalkulator \textit{Shift-Reduce} dalam C}
  \begin{columns}[T,onlytextwidth]
    \begin{column}{0.58\textwidth}
      \lstinputlisting[language=C,firstline=30,lastline=38,
        firstnumber=30,
        caption={Fungsi \texttt{reduce\_top} pada
        \file{shift\_reduce.c}},
        label={lst:p10-reduce}]{../buku_ajar/code/teaching/shift_reduce.c}
    \end{column}
    \begin{column}{0.38\textwidth}
      \begin{itemize}
        \item Pop dua nilai dan satu operator.
        \item Hitung hasil dengan \texttt{apply}.
        \item Dorong hasil kembali.
        \item Program mencetak urutan shift dan reduce.
        \item Masukan \texttt{2+3*4} menghasilkan 14.
      \end{itemize}
    \end{column}
  \end{columns}
\end{frame}

% ---------------------------------------------------------------------
\begin{frame}{Deteksi Kesalahan}
  \begin{block}{Kapan kesalahan terdeteksi}
    \small
    Ketika relasi precedence tidak terdefinisi untuk pasangan terminal
    yang berdekatan.
  \end{block}

  \begin{exampleblock}{Contoh}
    \footnotesize
    Pada \(\code{id}++\code{id}\), setelah shift \(+\) kedua, tabel
    tidak memiliki relasi antara \(+\) dan \(+\). Parser melaporkan
    kesalahan pada token kedua dan dapat melewati token sampai titik
    aman.
  \end{exampleblock}

  \vspace{2pt}
  {\footnotesize\color{ubbblue!80}Operator precedence tidak menangani
  grammar dengan produksi \(\varepsilon\) atau dua nonterminal
  bersebelahan; gunakan parser LR untuk grammar yang lebih umum.}
\end{frame}

% ---------------------------------------------------------------------
\begin{frame}{Perbandingan Top-Down dan Bottom-Up}
  \scriptsize
  \setlength{\tabcolsep}{4pt}
  \begin{tabular}{@{}p{2.6cm}p{5.2cm}p{5.8cm}@{}}
    \toprule
    \textbf{Aspek} & \textbf{Top-down} & \textbf{Bottom-up} \\
    \midrule
    Arah & Akar ke daun & Daun ke akar \\
    \textit{Left recursion} & Harus dihapus & Dapat ditangani \\
    Grammar & LL(1) dan variannya & LR, SLR, LALR, precedence \\
    Tabel & Satu entri per nonterminal & Satu entri per state \\
    Pesan kesalahan & Mudah kontekstual & Lebih sulit \\
    Implementasi & Recursive descent & Stack dan tabel \\
    Contoh & Parser LL(1) & Bison, Yacc, operator precedence \\
    \bottomrule
  \end{tabular}

  \vspace{4pt}
  {\scriptsize\color{ubbblue!75}Operator precedence adalah bentuk
  sederhana bottom-up; LR pada Pertemuan 11 lebih umum.}
\end{frame}

% ---------------------------------------------------------------------
\begin{frame}{Kegiatan Mahasiswa di Kelas}
  \begin{block}{Lokakarya: penelusuran dan kalkulator}
    \footnotesize
    \begin{enumerate}
      \setlength{\itemsep}{1pt}
      \item Telusuri shift-reduce untuk
            \(\code{id}*\code{id}+\code{id}\).
      \item Tandai handle pada setiap langkah reduce.
      \item Bangun tabel operator precedence untuk
            \(+\), \(-\), \(*\), \(/\), dan kurung.
      \item Modifikasi \file{shift\_reduce.c} agar mendukung
            \(-\), \(/\), dan kurung.
      \item Uji dengan ekspresi yang mengandung precedence dan
            associativity berbeda.
    \end{enumerate}
  \end{block}

  \vspace{4pt}
  {\footnotesize\color{ubbblue!80}Minimal 10 ekspresi uji. Data
  penelusuran dan tabel precedence disimpan untuk Pertemuan 11.}
\end{frame}

% ---------------------------------------------------------------------
\begin{frame}{Tugas dan Artefak}
  \begin{block}{Artefak Sub-CPMK 4.2 (bagian pertama)}
    Tabel operator precedence, penelusuran \textit{shift-reduce},
    kalkulator C, dan catatan konflik \textit{shift-reduce} atau
    \textit{reduce-reduce} yang ditemukan.
  \end{block}

  \small
  \begin{itemize}
    \item Penelusuran manual benar dan setiap handle ditandai.
    \item Tabel precedence lengkap untuk empat operator dan kurung.
    \item Kalkulator mendukung \(+\), \(-\), \(*\), \(/\), dan kurung.
    \item Minimal 10 ekspresi uji dengan hasil yang benar.
    \item Konflik yang muncul dicatat untuk \eng{debug} tabel
          ACTION/GOTO pada Pertemuan 11.
  \end{itemize}

  \vspace{3pt}
  {\footnotesize\color{ubbblue!75}Artefak lengkap Sub-CPMK 4.2
  (tabel aksi dan \eng{goto} SLR/LALR) diselesaikan pada Pertemuan 11
  dan dilaporkan pada Praktikum 3.}
\end{frame}

% ---------------------------------------------------------------------
\begin{frame}{Rangkuman dan Jembatan}
  \begin{itemize}
    \item Parser bottom-up membangun pohon dari daun ke akar dengan
          \eng{shift} dan \eng{reduce}.
    \item Handle adalah substring yang dapat direduksi; viable prefix
          menentukan stack yang masih valid.
    \item Konflik \eng{shift-reduce} dan \eng{reduce-reduce}
          diselesaikan dengan mengubah grammar atau kebijakan tetap.
    \item Operator precedence memakai relasi \(<\cdot\), \(\doteq\),
          \(\cdot>\) dan tabel yang lebih kecil daripada tabel LR.
    \item Data penelusuran dan tabel precedence menjadi fondasi tabel
          ACTION/GOTO SLR/LALR.
  \end{itemize}

  \vspace{4pt}
  \begin{exampleblock}{Pertemuan berikutnya}
    Pertemuan 11 membangun item LR(0), tabel ACTION/GOTO SLR, canonical
    LR, dan LALR, lalu men-\eng{debug} konflik yang muncul.
  \end{exampleblock}
\end{frame}

\end{document}
